\section{The objective gap and the reduced Hessian}
\label{sec:setup}

In the main geometric results, $P$ is a compact convex subset of a
finite-dimensional Euclidean space, with affine hull $x^0+\V$ and relative
interior $P^\circ=\ri P$. The tangent space $\V$ has dimension $d\geq1$.
We minimize a linear functional $\inner{c}{x}$ that is nonconstant on $P$;
equivalently, its orthogonal projection $c_{\V}$ onto $\V$ is nonzero. Set
\begin{align}
 \OPT&=\min_{x\in P}\inner{c}{x}, &
 \Phi&=\{x\in P:\inner{c}{x}=\OPT\},\label{eq:optimum}\\
 g(x)&=\inner{c}{x}-\OPT, &
 \Delta&=\max_{x\in P}g(x)>0.\label{eq:gap-range}
\end{align}
For $0\leq g\leq\Delta$, define
\begin{equation}
 L(g)=\{y\in P:g(y)\leq g\},\qquad
 D(g)=\diam L(g),\qquad K_g=L(g)-L(g)\subset\V.
 \label{eq:sublevels}
\end{equation}
The difference body $K_g$ consists of all chords of the sublevel set,
translated to the origin. Unlike the diameter, it records the longest
chord in every direction. For an interior point $x$, also put
\begin{equation}
 \ell(x)=\max_{y\in L(g(x))}\norm{y-x},\qquad
 \frac{D(g(x))}{2}\leq\ell(x)\leq D(g(x)).
 \label{eq:chord-length}
\end{equation}
The lower inequality follows from the triangle inequality applied to a
pair of points at distance $D(g(x))$.

\begin{definition}[Barrier and metric]
\label{def:barrier}
A $\nu$-self-concordant barrier for $P$, with $\nu>0$, is a convex
$C^3$ function $F:P^\circ\to\R$ that tends to $+\infty$ at every relative
boundary point, has positive-definite tangent Hessian, and satisfies
\begin{align}
 |D^3F(x)[h,h,h]|&\leq2\bigl(D^2F(x)[h,h]\bigr)^{3/2},\label{eq:sc}\\
 |DF(x)[h]|&\leq\sqrt\nu\bigl(D^2F(x)[h,h]\bigr)^{1/2}
 \qquad(h\in\V).\label{eq:barrier-gradient}
\end{align}
All gradients and Hessians below are taken along the affine space
$x^0+\V$, that is, in directions $h\in\V$. Write $H_F(x)$ for the
self-adjoint operator on $\V$ representing $D^2F(x)$; we call it the
reduced Hessian. Define
\begin{equation}
 \norm{h}_x=\sqrt{\inner{H_F(x)h}{h}},\qquad
 \norm{v}_x^*=\sqrt{\inner{H_F(x)^{-1}v}{v}},\qquad
 E_F(x)=\{h\in\V:\norm{h}_x\leq1\}.
 \label{eq:local-metric}
\end{equation}
The condition number is
$\kappa_F(x)=\lambda_{\max}(H_F(x))/\lambda_{\min}(H_F(x))$.
\end{definition}

If $W$ is an orthonormal basis matrix for $\V$, the matrix of this operator
is $W^\top\nabla^2F(x)W$; its eigenvalues do not depend on the choice of
orthonormal basis. The Euclidean metric is fixed throughout. For symmetric
matrix variables it is the Frobenius metric, and for several blocks it is
the product Frobenius metric. We use the Hessian of $F$, not of $\mu F$.
Multiplying by $\mu$ leaves the condition number unchanged but multiplies
every eigenvalue by $\mu$. A different representation of the feasible set,
or a nonorthogonal change of coordinates, needs a separate analysis of
the metric.

\subsection{Two standard barrier facts, with proofs}

The following facts are standard properties of self-concordant barriers
\citep{nesterov1994,renegar2001}. We give short one-dimensional proofs
to fix the normalization and the treatment of boundary points.

\begin{lemma}[Dikin containment and semiboundedness]
\label{lem:standard-containment}
For $x\in P^\circ$,
\begin{equation}
 x+\operatorname{int}_{\V} E_F(x)\subset P^\circ,
 \qquad x+E_F(x)\subset P,
 \label{eq:dikin}
\end{equation}
and for every $y\in P$,
\begin{equation}
 \inner{\nabla F(x)}{y-x}\leq\nu.
 \label{eq:semiboundedness}
\end{equation}
In particular, at every interior point, whether central or not,
\begin{equation}
 \norm{c_{\V}}_x^*\leq g(x).
 \label{eq:objective-support}
\end{equation}
\end{lemma}

\begin{proof}
Fix $h\in\V$ with $\norm{h}_x=1$, and write
$\varphi(t)=F(x+th)$ for $0\leq t<R$, where $R$ is the first positive
exit time from $P^\circ$. Self-concordance gives
\begin{equation}
 \left|\frac{d}{dt}\frac1{\sqrt{\varphi''(t)}}\right|\leq1.
 \label{eq:line-hessian}
\end{equation}
Thus $\varphi''(t)\leq(1-t)^{-2}$ for $0\leq t<\min\{R,1\}$.
If $R<1$, integration twice bounds $\varphi(t)$ as $t\uparrow R$,
contradicting boundary divergence. Hence $R\geq1$. This proves the open
containment, and closed containment follows by taking limits in the closed
set $P$.

To prove \eqref{eq:semiboundedness}, note that the case $y=x$ is
immediate. Otherwise let $\psi(t)=F(x+t(y-x))$, $0\leq t<1$. These points
are interior because $P$ is convex and $x$ is in the relative interior.
If $\psi'(0)\leq0$, the claim is immediate. If $\psi'(0)>0$,
convexity and \eqref{eq:barrier-gradient} give
\[
 \psi'(t)>0,\qquad
 \psi''(t)\geq\frac{\psi'(t)^2}{\nu},\qquad
 \frac1{\psi'(s)}\leq\frac1{\psi'(0)}-\frac{s}{\nu}
 \quad(0\leq s<1).
\]
The left side is positive. Letting $s\uparrow1$ proves
$\psi'(0)\leq\nu$, including boundary choices of $y$.
Finally, minimizing the linear objective over the closed Dikin ellipsoid
gives $\inner{c}{x}-\norm{c_{\V}}_x^*\geq\OPT$.
\end{proof}

\subsection{Existence and objective-gap parameterization}

For $\mu>0$, the central point is
\begin{equation}
 x_F(\mu)=\argmin_{x\in P^\circ}
       \left(F(x)+\frac{\inner{c}{x}}{\mu}\right),
 \qquad \nabla F(x_F(\mu))=-\frac{c_{\V}}{\mu}.
 \label{eq:central-path}
\end{equation}
We call $g_F(\mu)=g(x_F(\mu))$ the \emph{primal objective gap} of the
central point. Even for the logarithmic barrier, it need not equal the
primal--dual gap $\nu\mu$.

Parameterizing central paths by the objective value is classical.
In particular, the bounds in the following proposition agree with
\citet[Proposition 3.2]{pena2002}. Our proof uses only the Dikin bound
\eqref{eq:objective-support} and differentiation of the centrality
condition.

\begin{proposition}[The attained gap interval]
\label{prop:gap-parametrization}
The analytic center $x_F^{\rm ac}=\argmin F$ and all points
$x_F(\mu)$ exist and are unique. Put $a_F=g(x_F^{\rm ac})\in(0,\Delta)$.
Then $g_F$ is continuously differentiable, strictly increasing, and maps
$(0,\infty)$ bijectively onto $(0,a_F)$. Moreover,
\begin{align}
 g_F'(\mu)&=\frac{\norm{c_{\V}}_{x_F(\mu)}^{*\,2}}{\mu^2}>0,
 \label{eq:gap-derivative}\\
 \frac{\mu a_F}{\mu+a_F}&\leq g_F(\mu)\leq\nu\mu.
 \label{eq:gap-mu-comparison}
\end{align}
Thus $g_F(\mu)=\Theta(\mu)$ as $\mu\downarrow0$ for every fixed barrier.
\end{proposition}

\begin{proof}
A supporting affine function of $F$ bounds it below on the compact set
$P$. Since $F$ diverges at the boundary, its bounded sublevel sets are
compact subsets of $P^\circ$. The same holds after adding any linear function. Existence
follows, and the positive-definite Hessian gives strict convexity and
uniqueness. The implicit function theorem applied to
\eqref{eq:central-path} gives
$x_F'(\mu)=H_F(x_F(\mu))^{-1}c_{\V}/\mu^2$, proving
\eqref{eq:gap-derivative}.

For any $x^*\in\Phi$, centrality and semiboundedness give
\[
 g_F(\mu)=\mu\inner{\nabla F(x_F(\mu))}{x^*-x_F(\mu)}
 \leq\nu\mu.
\]
Hence $g_F(\mu)\to0$ as $\mu\downarrow0$. For $\mu\to\infty$,
optimality gives $F(x_F(\mu))\leq F(x_F^{\rm ac})+\Delta/\mu$.
Since a fixed sublevel set of $F$ is compact and $F$ has a unique
minimizer, $x_F(\mu)\to x_F^{\rm ac}$. An interior point cannot minimize or
maximize a nonconstant linear functional, so $0<a_F<\Delta$.
Strict monotonicity proves the claimed range.

By \eqref{eq:objective-support} and \eqref{eq:gap-derivative},
$(1/g_F)'\geq-1/\mu^2$. Integration from $\mu$ to $M$ followed by
$M\to\infty$ yields
$1/g_F(\mu)\leq1/a_F+1/\mu$, which is the remaining inequality.
\end{proof}

Let $\mu_F:(0,a_F)\to(0,\infty)$ be the inverse of $g_F$.
We write the path parameterized by the gap, and its Hessian and
condition number, as
\begin{equation}
 \widehat x_F(g)=x_F(\mu_F(g)),\qquad
 \widehat H_F(g)=H_F(\widehat x_F(g)),\qquad
 \widehat\kappa_F(g)=\kappa_F(\widehat x_F(g)).
 \label{eq:gap-path-notation}
\end{equation}
Two fixed barriers can therefore be compared at equal gap on the common
interval $(0,\min\{a_F,a_G\})$. Equal values of $g$ generally correspond
to different values of $\mu$. This matters when the Hessians themselves,
and not only their condition numbers, are compared.
